% TOP_PROBLEM: {"problemId": "problem.polynomial-time-distribution-free-pac-learning-of-dnf-formulas", "canonicalTitle": "Polynomial-Time Distribution-Free PAC Learning of DNF Formulas", "releaseRank": 219, "primaryDomain": "theoretical_computer_science", "primaryDomainLabel": "Theoretical computer science", "displayStatus": "open_with_solved_subcases", "releaseStatus": "open_with_solved_subcases"}
% !TeX program = lualatex
% Definition and sources only; this file is not an LLM solution attempt.
\documentclass[11pt]{article}
\usepackage[margin=25mm]{geometry}
\usepackage{fontspec}
\setmainfont{Latin Modern Roman}
\usepackage{amsmath,amssymb,mathtools}
\usepackage{unicode-math}
\setmathfont{Latin Modern Math}
\usepackage{xurl,hyperref}
\hypersetup{colorlinks=true,urlcolor=blue}
\setcounter{secnumdepth}{0}
\setlength{\parindent}{0pt}
\setlength{\parskip}{5pt}
\setlength{\emergencystretch}{3em}
\begin{document}
\section*{\texorpdfstring{Polynomial-Time Distribution-Free PAC Learning of DNF Formulas}{Polynomial-Time Distribution-Free PAC Learning of DNF Formulas}}\label{219}
\subsection{Definitions and mathematical statement}
An $s$-term DNF on $\{0,1\}^n$ is an OR of at most $s$ terms, each an AND of literals $x_i$ or $\neg x_i$. The learner receives independent labeled examples $(x,f(x))$ with $x\sim D$, where $D$ is arbitrary and unknown. For $0<{\varepsilon},\delta<1$, it must output an efficiently evaluable hypothesis $h$ such that
\[
 \Pr_{\mathrm{samples,coins}}\bigl[\Pr_{x\sim D}(h(x)\ne f(x))\le{\varepsilon}\bigr]\ge1-\delta
\]
using time polynomial in $n,s,1/{\varepsilon},\log(1/\delta)$. The hypothesis need not itself be a DNF. Membership queries or a known special input distribution are not supplied by the target.
\par\smallskip\textit{Formulation and concept references:} \href{https://arxiv.org/html/2505.18839v1}{[S1]}.

\subsection{Short English statement}
Can arbitrary DNF formulas be learned efficiently from random labeled examples drawn from any unknown distribution, even when the learner may output a different kind of hypothesis?
\subsection{Sources}
\begin{itemize}
\item[S1]Alman, Nadimpalli, Patel, and Servedio, DNF Learning via Locally Mixing Random Walks, Section 1.1.
\url{https://arxiv.org/html/2505.18839v1}.
\textit{Formulation source.}
\end{itemize}
\end{document}
